\documentclass[10pt,a4paper]{amsart}
\usepackage[margin=19mm]{geometry}
\usepackage{amsmath,amssymb,mathtools}
\usepackage[scr=rsfs,cal=euler]{mathalfa}
\usepackage{xfrac}
\usepackage[unicode,hidelinks,bookmarks=false]{hyperref}
\newcommand{\ie}{i.e.\ }
\newcommand{\cf}{cf.\ }
\newcommand{\resp}{respectively\ }
\newcommand{\Subsection}{\subsection}
\providecommand{\nobreakdash}{\nobreak\hbox{-}}
\setlength{\emergencystretch}{2em}
\multlinegap0pt
\usepackage{bm}
\usepackage{bbm}
\usepackage{dsfont}
\usepackage{xspace}
\usepackage{cancel}
\usepackage{mathtools}
\usepackage{amsthm}
\makeatletter
\@ifundefined{theorem}{\newtheorem{theorem}{Theorem}}{}
\@ifundefined{assumption}{\newtheorem{assumption}[theorem]{Assumption}}{}
\@ifundefined{lemma}{\newtheorem{lemma}[theorem]{Lemma}}{}
\@ifundefined{proposition}{\newtheorem{proposition}[theorem]{Proposition}}{}
\makeatother
\newcommand{\obsdata}{\mathsf{d}}
\renewcommand{\d}{\mathrm{d}}
\title[Measure flow path recovery in Bayes Hilbert spaces]{Measure flow path recovery in Bayes Hilbert spaces}
\author{S. David Mis, Maarten V. de Hoop}
\date{Source: arXiv:2603.20329v2 (CC BY 4.0)}
\begin{document}
\maketitle

\noindent\emph{Source and licence.} Measure flow path recovery in Bayes Hilbert spaces. Version arXiv:2603.20329v2 (CC BY 4.0), distributed under the Creative Commons Attribution 4.0 International licence, as recorded in the arXiv licence metadata for this exact version and shown on the abstract page https://arxiv.org/abs/2603.20329v2. Full rights notice: licenses/arxiv-2603.20329-CC-BY-4.0.txt.\par\smallskip

\noindent\textit{Editorial scope.} Counted unit: the Theorem whose source label is \texttt{thm:pathwise-stability} in the article (source0.tex lines 2436--2468, source label \texttt{thm:pathwise-stability}), delivered together with the article's own complete original proof of it (source0.tex lines 2470--2673, one \texttt{proof} environment of 204 lines). No mathematics is written, repaired or re-derived: statement and proof are the article's own text line for line. Required context retained uncounted: ass:joint-transport-observability (lines 1789--1803); prop:transport-term-upper-bound (lines 1825--1839); ass:observation-regularity (lines 2351--2373); eq:J-lipschitz (lines 2366--2372); lem:a-priori-bounds (lines 2375--2418); eq:a-priori-temporal (lines 2408--2417). Every definition, display and label of those lines is retained unchanged. Omitted from this package and not redistributed: 1--1788, 1804--1824, 1840--2350, 2373--2374, 2418--2435, 2469--2469, 2674--5286. Nothing inside a retained excerpt is omitted or altered: every retained body is delivered exactly as the article's lines are written, so no omission receipt of any kind accompanies this package. Citations: the retained excerpts carry no citation command at all, so no bibliography entry of the article is transcribed and no reference is redistributed. Reference adaptation: none was needed and none was made; every reference inside the delivered unit resolves inside it, so no unresolved reference or citation is produced. Printed slips kept literal: no printed word, symbol, hypothesis or number of the counted statement or of its proof was altered. Layout and class: the article's own class is \texttt{article} with packages including natbib, inputenc, amsmath, amssymb, amsthm, xcolor, geometry, hyperref; the wrapper is the lane's pinned \texttt{amsart} editorial wrapper at 10pt on A4, which restates the theorem-like environments the retained text needs and adds the article's own macro definitions that the retained text needs (\texttt{\textbackslash d}, \texttt{\textbackslash obsdata}), with extra wrapper packages (bm, bbm, dsfont, xspace, cancel, mathtools). The delivered numbering is the unit's own. Assets: no retained span contains a figure environment, a graphics-inclusion command, a PDF-page inclusion, a table or a TikZ picture; no image file is copied into this package, the package declares no asset dependency and no image file is redistributed with it. Licence: version arXiv:2603.20329v2 of this article is distributed under the Creative Commons Attribution 4.0 International licence, as recorded in the arXiv licence metadata for this exact version and shown on the abstract page https://arxiv.org/abs/2603.20329v2. A full rights notice is delivered at licenses/arxiv-2603.20329-CC-BY-4.0.txt. Characters: every retained excerpt is pure ASCII, so the delivered engine, XeLaTeX with its default Latin Modern text font, reports no missing character.
\begin{assumption}[Joint transport--observability coercivity]
    \label{ass:joint-transport-observability}
    There exists \(\kappa>0\) such that for all \(h\in \mathcal A_{\mathrm{ad}}\) and all
    \[
        \xi\in L^2(0,T;H^{s'}(\Omega)\cap L^2_0(\nu_0))
        \cap H^1(0,T;L^2_0(\nu_0)),
    \]
    we have
    \begin{equation}
        \label{eq:joint-coercivity}
        Q_y[h](\xi,\xi)
        \ge
        \kappa \|\xi\|_{L^2(0,T;L^2(\nu_0))}^2.
    \end{equation}
\end{assumption}

\begin{proposition}[Upper bound on the transport term]
    \label{prop:transport-term-upper-bound}
    There exists a constant \(C_{\mathrm{tr}}>0\), depending only on the admissible class
    \(\mathcal X_{\mathrm{ad}}\) and the domain \(\Omega\), such that for all
    \(h\in\mathcal A_{\mathrm{ad}}\) and all
    \[
        \xi\in H^1(0,T;L^2_0(\nu_0)),
    \]
    we have
    \[
        \int_0^T \mathfrak g_{h(t)}(\dot\xi(t),\dot\xi(t))\,\d t
        \le
        C_{\mathrm{tr}}\|\dot\xi\|_{L^2(0,T;L^2(\nu_0))}^2.
    \]
\end{proposition}

\begin{assumption}[Regularity of the mobile observation family]
    \label{ass:observation-regularity}
    Assume that there exists an open neighborhood \(U\subset \mathcal Z\) of
    \(\mathcal X_{\mathrm{ad}}\), together with constants \(r_0>0\) and \(L>0\), such that for each
    \(t\in[0,T]\),
    \[
        \mathcal G_t^y:U\to\mathbb R^r
    \]
    is \(C^1\) with respect to the \(H^{s'}(\Omega)\)-topology, and its differential
    \[
        J_{t,h}^{y}=D\mathcal G_t^y(h)
    \]
    is locally Lipschitz in \(h\), uniformly in \(t\), as a map into
    \(\mathcal L(L^2(\nu_0),\mathbb R^r)\): for all \(h,k\in U\) with
    \(\|h-k\|_{H^{s'}(\Omega)}\le r_0\),
    \begin{equation}
        \label{eq:J-lipschitz}
        \|J_{t,h}^{y}-J_{t,k}^{y}\|_{\mathcal L(L^2(\nu_0),\mathbb R^r)}
        \le
        L\|h-k\|_{H^{s'}(\Omega)}
        \qquad\text{for all } t\in[0,T].
    \end{equation}
\end{assumption}

\begin{lemma}[A priori bounds from minimality]
    \label{lem:a-priori-bounds}
    Let \(h^\dagger\in\mathcal A_{\mathrm{ad}}\), and let
    \[
        \obsdata(t)=\mathcal G_t^y(h^\dagger(t))+\eta(t)
    \]
    with
    \[
        \|\eta\|_{L^2(0,T;\mathbb R^r)}\le \delta.
    \]
    Define the regularization energy of the true path by
    \begin{equation}
        \label{eq:regularization-energy-true-path}
        R[h^\dagger]
        :=
        \frac{\lambda}{2}\int_0^T \mathfrak g_{h^\dagger(t)}(\dot h^\dagger(t),\dot h^\dagger(t))\,\d t
        +
        \frac{\mu}{2}\int_0^T
        \Big(
        \|h^\dagger(t)\|_{L^2(\nu_0)}^2+\|\dot h^\dagger(t)\|_{L^2(\nu_0)}^2
        \Big)\,\d t
        +
        \frac{\gamma}{2}\int_0^T \|h^\dagger(t)\|_{H^s(\Omega)}^2\,\d t.
    \end{equation}
    If \(h\in\mathcal A_{\mathrm{ad}}\) is a minimizer of \(\mathcal I_{\lambda,\mu,\gamma}^{\,y}\) for
    data \(\obsdata\), then
    \begin{equation}
        \label{eq:a-priori-functional-bound}
        \mathcal I_{\lambda,\mu,\gamma}^{\,y}[h]
        \le
        \tfrac12\delta^2+R[h^\dagger].
    \end{equation}
    In particular,
    \begin{align}
        \label{eq:a-priori-data-misfit}
        \tfrac12\|\mathcal G_\cdot^y(h(\cdot))-d\|_{L^2(0,T;\mathbb R^r)}^2
         & \le
        \tfrac12\delta^2+R[h^\dagger], \\
        \label{eq:a-priori-temporal}
        \tfrac{\mu}{2}\|\dot h\|_{L^2(0,T;L^2(\nu_0))}^2
         & \le
        \tfrac12\delta^2+R[h^\dagger].
    \end{align}
\end{lemma}

\begin{theorem}[Local pathwise stability under partial observability]
    \label{thm:pathwise-stability}
    Assume Assumptions~\ref{ass:joint-transport-observability}
    and~\ref{ass:observation-regularity}. Let
    \(h^\dagger\in \mathcal A_{\mathrm{ad}}^{s'}\), and let
    \[
        \obsdata(t)=\mathcal G_t^y(h^\dagger(t))+\eta(t)
    \]
    with
    \[
        \|\eta\|_{L^2(0,T;\mathbb R^r)}\le\delta.
    \]
    Let \(h\in\mathcal A_{\mathrm{ad}}^{s'}\) be a minimizer
    of \(\mathcal I_{\lambda,\mu,\gamma}^{\,y}\) for data \(\obsdata\).

    Then there exist constants \(r>0\) and \(C>0\), depending only on
    \(\mathcal X_{\mathrm{ad}}\), \(\Omega\), \(y\), the coercivity constant \(\kappa\), and
    the regularity constants in Assumption~\ref{ass:observation-regularity}, such that if
    \begin{equation}
        \label{eq:proximity-condition}
        \sup_{t\in[0,T]}\|h(t)-h^\dagger(t)\|_{H^{s'}(\Omega)}\le r,
    \end{equation}
    then
    \begin{equation}
        \label{eq:pathwise-stability-estimate}
        \|h-h^\dagger\|_{L^2(0,T;L^2(\nu_0))}^2
        \le
        \frac{C}{\kappa}(1+\mu^{-1})
        \Big(
        \delta^2+R[h^\dagger]
        \Big).
    \end{equation}
\end{theorem}

\begin{proof}
    Set \(\xi:=h-h^\dagger\). We estimate \(\|\xi\|_{L^2(0,T;L^2(\nu_0))}^2\) by applying the joint
    coercivity of \(Q_y[h^\dagger]\) and bounding the resulting terms using minimality and the
    proximity condition.

    \medskip
    \noindent
    \emph{Step 1: joint coercivity.}
    By Assumption~\ref{ass:joint-transport-observability} applied at the path \(h^\dagger\),
    \begin{equation}
        \label{eq:stab-joint-coercivity}
        \kappa\|\xi\|_{L^2_tL^2_x}^2
        \le
        \underbrace{
        \int_0^T \|J_{t,h^\dagger(t)}^{y}\xi(t)\|_{\mathbb R^r}^2\,\d t
        }_{\mathrm{Term\;I}}
        +
        \underbrace{
            \int_0^T \mathfrak g_{h^\dagger(t)}(\dot\xi(t),\dot\xi(t))\,\d t
        }_{\mathrm{Term\;II}}.
    \end{equation}

    \medskip
    \noindent
    \emph{Step 2: control of Term I.}
    Since \(h^\dagger\in \mathcal A_{\mathrm{ad}}^{s'}\), the image
    \(h^\dagger([0,T])\subset \mathcal Z\) is compact. Because \(U\subset\mathcal Z\) is an open
    neighborhood of \(\mathcal X_{\mathrm{ad}}\) and \(h^\dagger(t)\in\mathcal X_{\mathrm{ad}}\subset U\)
    for all \(t\), there exists \(r_U>0\) such that
    \[
        \bigcup_{t\in[0,T]} B_{\mathcal Z}(h^\dagger(t),r_U)\subset U.
    \]
    We choose \(r\) so that
    \[
        r\le \min\left\{r_0,\;r_U,\;\sqrt{\frac{\kappa}{L^2}}\right\}.
    \]
    Then for every \(t\in[0,T]\) and every \(\theta\in[0,1]\),
    \[
        h^\dagger(t)+\theta\xi(t)\in U,
    \]
    so the fundamental-theorem-of-calculus expansion and the Lipschitz estimate
    \eqref{eq:J-lipschitz} may be applied along the segment between
    \(h^\dagger(t)\) and \(h(t)\).

    By the fundamental theorem of calculus in Banach spaces,
    \[
        \mathcal G_t^y(h(t))-\mathcal G_t^y(h^\dagger(t))
        =
        J_{t,h^\dagger(t)}^{y}\xi(t)
        +
        \int_0^1
        \bigl(J_{t,h^\dagger(t)+\theta\xi(t)}^{y}-J_{t,h^\dagger(t)}^{y}\bigr)\xi(t)\,\d\theta.
    \]
    Denote the remainder by
    \[
        r(t)
        :=
        \int_0^1
        \bigl(J_{t,h^\dagger(t)+\theta\xi(t)}^{y}-J_{t,h^\dagger(t)}^{y}\bigr)\xi(t)\,\d\theta.
    \]
    By the Lipschitz condition \eqref{eq:J-lipschitz} and the proximity condition
    \eqref{eq:proximity-condition},
    \[
        \|r(t)\|_{\mathbb R^r}
        \le
        \int_0^1
        L\theta\|\xi(t)\|_{H^{s'}(\Omega)}\|\xi(t)\|_{L^2(\nu_0)}\,\d\theta
        \le
        \frac{Lr}{2}\|\xi(t)\|_{L^2(\nu_0)}.
    \]
    Since
    \[
        J_{t,h^\dagger(t)}^{y}\xi(t)
        =
        \mathcal G_t^y(h(t))-\mathcal G_t^y(h^\dagger(t))-r(t),
    \]
    we obtain
    \[
        \|J_{t,h^\dagger(t)}^{y}\xi(t)\|_{\mathbb R^r}^2
        \le
        2\|\mathcal G_t^y(h(t))-\mathcal G_t^y(h^\dagger(t))\|_{\mathbb R^r}^2
        +
        2\|r(t)\|_{\mathbb R^r}^2.
    \]
    For the first piece, write
    \[
        \mathcal G_t^y(h(t))-\mathcal G_t^y(h^\dagger(t))
        =
        \bigl(\mathcal G_t^y(h(t))-\obsdata(t)\bigr)+\eta(t),
    \]
    so that
    \[
        \|\mathcal G_t^y(h(t))-\mathcal G_t^y(h^\dagger(t))\|_{\mathbb R^r}^2
        \le
        2\|\mathcal G_t^y(h(t))-\obsdata(t)\|_{\mathbb R^r}^2+2\|\eta(t)\|_{\mathbb R^r}^2.
    \]
    Integrating in time and applying Lemma~\ref{lem:a-priori-bounds},
    \[
        \int_0^T \|\mathcal G_t^y(h(t))-\mathcal G_t^y(h^\dagger(t))\|_{\mathbb R^r}^2\,\d t
        \le
        4\bigl(\tfrac12\delta^2+R[h^\dagger]\bigr)+2\delta^2
        \le
        C_1\bigl(\delta^2+R[h^\dagger]\bigr).
    \]
    For the remainder,
    \[
        \int_0^T \|r(t)\|_{\mathbb R^r}^2\,\d t
        \le
        \frac{L^2 r^2}{4}\|\xi\|_{L^2_tL^2_x}^2.
    \]
    Therefore,
    \begin{equation}
        \label{eq:stab-term-I-bound}
        \mathrm{Term\;I}
        \le
        C_1\bigl(\delta^2+R[h^\dagger]\bigr)
        +
        \frac{L^2 r^2}{2}\|\xi\|_{L^2_tL^2_x}^2.
    \end{equation}

    \medskip
    \noindent
    \emph{Step 3: control of Term II.}
    By Proposition~\ref{prop:transport-term-upper-bound},
    \[
        \mathrm{Term\;II}
        =
        \int_0^T \mathfrak g_{h^\dagger(t)}(\dot\xi(t),\dot\xi(t))\,\d t
        \le
        C_{\mathrm{tr}}\|\dot\xi\|_{L^2_tL^2_x}^2.
    \]
    Since \(\xi=h-h^\dagger\), we have
    \[
        \dot\xi=\dot h-\dot h^\dagger,
    \]
    and therefore
    \[
        \|\dot\xi\|_{L^2_tL^2_x}^2
        \le
        2\|\dot h\|_{L^2_tL^2_x}^2
        +
        2\|\dot h^\dagger\|_{L^2_tL^2_x}^2.
    \]
    By the a priori bound \eqref{eq:a-priori-temporal},
    \[
        \frac{\mu}{2}\|\dot h\|_{L^2_tL^2_x}^2
        \le
        \frac12\delta^2+R[h^\dagger],
    \]
    so
    \[
        \|\dot h\|_{L^2_tL^2_x}^2
        \le
        \frac{1}{\mu}\bigl(\delta^2+2R[h^\dagger]\bigr).
    \]
    Also, by the definition of \(R[h^\dagger]\),
    \[
        R[h^\dagger]
        \ge
        \frac{\mu}{2}\|\dot h^\dagger\|_{L^2_tL^2_x}^2,
    \]
    hence
    \[
        \|\dot h^\dagger\|_{L^2_tL^2_x}^2
        \le
        \frac{2}{\mu}R[h^\dagger].
    \]
    Combining these estimates gives
    \begin{equation}
        \label{eq:stab-term-II-bound}
        \mathrm{Term\;II}
        \le
        \frac{C_3}{\mu}\bigl(\delta^2+R[h^\dagger]\bigr)
    \end{equation}
    for a constant \(C_3>0\) depending only on \(C_{\mathrm{tr}}\).

    \medskip
    \noindent
    \emph{Step 4: combine and absorb.}
    Substituting \eqref{eq:stab-term-I-bound} and \eqref{eq:stab-term-II-bound} into
    \eqref{eq:stab-joint-coercivity} gives
    \[
        \kappa\|\xi\|_{L^2_tL^2_x}^2
        \le
        \frac{L^2 r^2}{2}\|\xi\|_{L^2_tL^2_x}^2
        +
        C_1\bigl(\delta^2+R[h^\dagger]\bigr)
        +
        \frac{C_3}{\mu}\bigl(\delta^2+R[h^\dagger]\bigr).
    \]
    Choose
    \[
        r
        \le
        \min\left\{r_0,\;r_U,\;\sqrt{\frac{\kappa}{L^2}}\right\},
    \]
    so that \(\frac{L^2 r^2}{2}\le \frac\kappa2\). Then
    \[
        \frac\kappa2\|\xi\|_{L^2_tL^2_x}^2
        \le
        C(1+\mu^{-1})\bigl(\delta^2+R[h^\dagger]\bigr),
    \]
    and dividing by \(\kappa/2\) gives \eqref{eq:pathwise-stability-estimate}.
\end{proof}

\end{document}
